\documentclass[11pt]{article}
\usepackage{amsmath,amsthm,amssymb}
\usepackage[square,numbers]{natbib}
\bibliographystyle{plainnat}
\usepackage{hyperref}
\usepackage{geometry}

\theoremstyle{plain}
\newtheorem{definition}{Definition}[section]
\theoremstyle{remark}
\newtheorem*{remark}{Remark}

\title{Encrypted Search as Trapdoor Computing}
\author{Alexander Towell}
\date{Draft --- \today}

\begin{document}
\maketitle

\begin{abstract}
Encrypted search is an application of trapdoor computing in which an
untrusted provider evaluates cipher maps over confidential queries and
indexes without learning the queries, the indexes, or the results.
We formalize encrypted search as an instance of the cipher map abstraction,
mapping the standard encrypted search vocabulary (information needs, search
agents, secure indexes, hidden queries) to the four-property framework
of trapdoor computing. This positions encrypted search as one application
domain among many --- trapdoor computing in general can approximate
arbitrary computable functions.
\end{abstract}

%===========================================================================
\section{Introduction}
%===========================================================================

An \emph{information retrieval} process begins when a \emph{search agent}
submits a query representing an \emph{information need} to an information
system, which returns relevant objects (e.g., documents).

\emph{Encrypted search} is information retrieval in which an untrusted
provider evaluates queries over confidential data without learning
the queries or the data. The provider sees only opaque bit strings
and total functions on those bit strings --- the defining characteristics
of a cipher map.

Encrypted search is one application of trapdoor computing. The general
framework supports any computable function $f : X \to Y$ via cipher maps
$\hat{f} : \{0,1\}^n \to \{0,1\}^n$, with four properties (totality,
representation uniformity, correctness, composability) parameterized by
$(\eta, \varepsilon, \mu, \delta)$. See the cipher map formalism
for precise definitions.

%===========================================================================
\section{Vocabulary Mapping}
%===========================================================================

The encrypted search literature uses domain-specific terminology.
Each term maps directly to the cipher map abstraction.

\begin{definition}[Search Agent (SA)]
The \emph{search agent} is the trusted machine $T$. It holds the
trapdoor $s$, the encoding function $\mathrm{enc}$, and the decoding
function $\mathrm{dec}$. It formulates information needs, encodes
queries, and decodes results.
\end{definition}

\begin{definition}[Encrypted Search Provider (ESP)]
The \emph{encrypted search provider} is the untrusted machine $U$.
It holds cipher maps and cipher values. It evaluates
$\hat{f}(c)$ for any $c \in \{0,1\}^n$ and returns results to $T$.
It cannot decode, cannot distinguish real queries from filler, and
cannot determine which function is being computed.
\end{definition}

\begin{definition}[Information Need]
An \emph{information need} is a latent function $f : X \to Y$ that the
search agent wishes to evaluate. For keyword search, $f = \mathbf{1}_A$
is the indicator function for a set $A$ of documents containing a keyword.
For richer retrieval, $f$ may be a relevance scoring function, a ranked
retrieval function, or any computable map.
\end{definition}

\begin{definition}[Hidden Query]
A \emph{hidden query} is a cipher value $\mathrm{enc}(x, k)$ ---
the encoding of a plaintext query $x$ under representation index $k$
and trapdoor $s$. The ESP evaluates $\hat{f}(\mathrm{enc}(x, k))$
without learning $x$.
\end{definition}

\begin{definition}[Secure Index]
A \emph{secure index} is a cipher map $\hat{f}$ stored on the ESP.
In the general case, a secure index is any cipher map --- not
necessarily a Bloom filter or any specific data structure.

The four properties of a cipher map specialize to encrypted search as:
\begin{enumerate}
    \item \textbf{Totality}: The ESP can evaluate the index on any
    bit string. There is no ``not found'' signal; filler queries
    produce output indistinguishable from real queries.

    \item \textbf{Representation uniformity}: Hidden queries are
    $\delta$-close to uniform. Frequency analysis of query traffic
    is bounded by $\delta$.

    \item \textbf{Correctness}: The search agent recovers the correct
    answer with probability $\geq 1 - \eta$. Nonzero $\eta$ provides
    plausible deniability: the ESP cannot distinguish a true positive
    from a false positive.

    \item \textbf{Composability}: Boolean combinations of queries
    (AND, OR, NOT) compose as cipher map chains with predictable
    error accumulation $\eta_{\text{total}} = 1 - \prod(1 - \eta_i)$.
\end{enumerate}
\end{definition}

\begin{remark}[Generality]
Traditional encrypted search constructions use Bloom filters (which
output observable Boolean values) or searchable encryption schemes
(which leak access patterns). The cipher map abstraction subsumes both:
a Bloom filter is a cipher map with $\eta = 0$, $K(x) = 1$
(no representation uniformity), and observable output. A secure index
as defined here generalizes this --- the output is an opaque cipher
value decoded only by $T$, and the construction may be any cipher map
(batch or online, membership or arbitrary function).
\end{remark}

%===========================================================================
\section{Scope}
%===========================================================================

Encrypted search is one application of trapdoor computing. The cipher map
abstraction applies to any setting where:
\begin{itemize}
    \item A trusted party wants to delegate computation to an untrusted party,
    \item The function, its domain, and its codomain are confidential,
    \item Approximate answers are acceptable (or desirable for privacy).
\end{itemize}

Other applications include encrypted databases, private set intersection,
confidential analytics, and any computation delegated to an untrusted
cloud. The formalism is the same; only the choice of latent function
$f$ and the parameter trade-offs $(\eta, \varepsilon, \mu, \delta, p)$
change.

% TODO: When the formalism paper is complete, cite it here.
% TODO: Flesh out with concrete constructions for keyword search,
%       ranked retrieval, Boolean query composition.
% TODO: Experimental validation against the formalism's parameter
%       predictions.

\end{document}
